% This file is generated from the corresponding Obsidian note.
% Source: Teoricos/GS - Teo1.md
\chapter{Diferenciabilidad y regla de la cadena}
\label{ch:teo-1}
\section{Diferenciabilidad en espacios normados}
\begin{bookdefinition}{Definición}
Sean $V$ y $W$ espacios vectoriales reales de dimensión finita, normados. Sea $I \subseteq V$ un abierto y $p \in I$. Una función $F \colon I \to W$ se dice \emph{diferenciable en $p$} si existe una transformación lineal $T \colon V \to W$ tal que

\[
\lim_{v \to 0} \frac{\|F(p+v)-F(p)-Tv\|}{\|v\|} = 0.
\]
La norma del denominador es la de $V$ y la del numerador es la de $W$.

\end{bookdefinition}
\begin{bookremark}{Pregunta}
¿Qué representa $Tv$?

\end{bookremark}
\begin{booklemma}{}
Sean $U$, $W$ y $X$ espacios vectoriales reales normados de dimensión finita. Sean

\[
g \colon A \subseteq U \to W,\qquad f \colon B \subseteq W \to X
\]
funciones con $\operatorname{Im}(g)\subseteq B$. Sean $x_0$ e $y_0$ puntos de acumulación de $A$ y $B$, respectivamente. Supongamos que

\[
\lim_{x\to x_0} g(x)=y_0, \qquad \lim_{y\to y_0} f(y)=z.
\]
Si se cumple alguna de las siguientes condiciones:

\begin{enumerate}
\item $f$ es continua en $y_0$.
\item $g(x)\neq y_0$ en una vecindad agujereada de $x_0$. Entonces
\end{enumerate}
\[
\lim_{x\to x_0} f(g(x))=z.
\]
\end{booklemma}
\begin{bookproposition}{Matriz Diferencial}
Sea $F \colon I \subseteq V \to W$ diferenciable en $p$ y sea $T \colon V \to W$ una transformación lineal tal que

\[
\lim_{v \to 0} \frac{\|F(p+v)-F(p)-Tv\|}{\|v\|} = 0.
\]
Entonces, para todo $v \in V$,

\[
Tv = \lim_{t \to 0} \frac{F(p+tv)-F(p)}{t}.
\]
\end{bookproposition}
\begin{bookproof}{}
Si $v=0$, la igualdad es inmediata porque $T0=0$. Supongamos $v \neq 0$. Como $I$ es abierto y $p \in I$, existe $\varepsilon>0$ tal que $p+tv \in I$ para todo $t \in (-\varepsilon,\varepsilon)$. Entonces

\[
0
= \lim_{t \to 0}
\frac{\|F(p+tv)-F(p)-T(tv)\|}{\|tv\|}
= \lim_{t \to 0}
\frac{1}{\|v\|}
\left\|\frac{F(p+tv)-F(p)}{t}-Tv\right\|.
\]
Multiplicando por $\|v\|$, se obtiene

\[
\lim_{t \to 0}\left\|\frac{F(p+tv)-F(p)}{t}-Tv\right\|=0,
\]
de donde

\[
Tv = \lim_{t \to 0} \frac{F(p+tv)-F(p)}{t}.
\]
\bookqed
\end{bookproof}
\begin{bookcorollary}{Corolario}
Si $F$ es diferenciable en $p$, existe una única transformación lineal que satisface la definición anterior. Se la llama \emph{diferencial de $F$ en $p$} y se la denota por

\[
(\mathrm{d}F)_p \colon V \to W.
\]
\end{bookcorollary}
\begin{bookproof}{}
La proposición anterior muestra que cualquier transformación lineal que satisfaga la definición debe verificar

\[
Tv = \lim_{t \to 0} \frac{F(p+tv)-F(p)}{t}
\qquad \text{para todo } v \in V,
\]
y por lo tanto queda completamente determinada por $F$ y por $p$.

\bookqed
\end{bookproof}
\begin{bookdefinition}{Aproximacion afin}
Sea $F \colon I \subseteq V \to W$ es diferenciable en $p$. Podemos definir

\[
\begin{aligned} R(v)&= F(p+v)-F(p)-Tv\\ & =  F(p+v)-F(p)-(dF)_{p}(v)\\&=F(p+v)-\bigg(F(p)+(dF)_{p}(v)\bigg)\end{aligned}
\]
Sustituyendo llegamos a que

\[
R(x-p)=F(x)-\bigg(F(p)+(dF)_{p}(x-p)\bigg)
\]
El término $R(v)$ es el error al aproximar $F(p+v)$ por la \emph{aproximación afín de $F$ en $p$}

\[
A_{p}(x)=F(p)+(\mathrm{d}F)_p(x-p).
\]
Osea podemos pensar

\[
F(x)\approx A_{p}(x)
\]
O podemos pensar

\[
F(x)=A_{p}(x)+R(x-p)
\]
donde $R(x-p)$ cumple $\lim\limits_{ x \to p }\lVert R(x-p) \rVert=0$. Mas aún, lo hace mas rapido que lineal

\[
\lim_{x\to p}\frac{\|R(x-p)\|}{\|x-p\|}=0.
\]
\end{bookdefinition}
\begin{booklemma}{Lema}
Toda transformación lineal $T \colon V \to W$ entre espacios vectoriales reales normados de dimensión finita es acotada. Es decir, existe $\mu \ge 0$ tal que

\[
\|Tv\|\le \mu \|v\|
\qquad \text{para todo } v \in V.
\]
\end{booklemma}
\begin{bookproof}{}
\begin{enumerate}
\item Sea
\end{enumerate}
\[
S=\{v\in V:\|v\|=1\}
\]
\begin{enumerate}[start=2]
\item La función $G \colon S \to \mathbb{R}$, $G(v)=\|Tv\|$, es continua.
\item Como $S$ es cerrado y acotado en un espacio de dimensión finita, es compacto; por ende $G(S)$ alcanza un máximo, digamos $\mu\ge 0$.
\item Luego
\end{enumerate}
\[
\|Tv\|\le \mu\qquad \text{para todo } v\in S.
\]
\begin{enumerate}[start=5]
\item Si $v\neq 0$, entonces $v/\|v\|\in S$, de modo que
\end{enumerate}
\[
\|Tv\|=\|v\|\,\left\|T\!\left(\frac{v}{\|v\|}\right)\right\|\le\mu \|v\|.
\]
\begin{enumerate}[start=6]
\item El caso $v=0$ es trivial.
\end{enumerate}
\bookqed
\end{bookproof}
\begin{booktheorem}{Teorema}
Si $F \colon I \subseteq V \to W$ es diferenciable en $p$, entonces $F$ es continua en $p$.

\end{booktheorem}
\begin{bookproof}{}
Sea $T=(\mathrm{d}F)_p$. Entonces

\[
F(p+v)-F(p)=\bigl(F(p+v)-F(p)-Tv\bigr)+Tv.
\]
Tomando normas y usando el lema anterior,

\[
\|F(p+v)-F(p)\|
\le
\|F(p+v)-F(p)-Tv\|+\|Tv\|
\le
\|v\|\,\frac{\|F(p+v)-F(p)-Tv\|}{\|v\|}+\mu \|v\|.
\]
Cuando $v\to 0$, ambos sumandos tienden a $0$, luego

\[
\lim_{v\to 0} F(p+v)=F(p).
\]
\bookqed
\end{bookproof}
\begin{bookexercise}{Ejercicio}
Rehacer la demostración del lema anterior usando continuidad de transformaciones lineales, o bien demostrar que toda transformación lineal entre espacios normados de dimensión finita es continua.

\end{bookexercise}
\section{Diferenciabilidad en conjuntos y regla de la cadena}
\begin{bookdefinition}{Definición}
Sea $F \colon I \subseteq V \to W$ con $I$ abierto. Se dice que $F$ es \emph{diferenciable} si es diferenciable en cada punto de $I$.

\end{bookdefinition}
\begin{bookdefinition}{Definición}
Sea $A\subseteq V$ arbitrario y sea $F \colon A \to W$ continua. Se dice que $F$ es \emph{diferenciable en $p\in A$} si existe un abierto $I_p\subseteq V$ con $p\in I_p$ y una función diferenciable

\[
F_p \colon I_p \to W
\]
tal que $F_p$ extiende a $F$ en $A\cap I_p$, es decir,

\[
F_p|_{A\cap I_p}=F|_{A\cap I_p}.
\]
\end{bookdefinition}
\begin{bookremark}{Observación}
Más adelante se puede probar una caracterización equivalente: $F \colon A \to W$ es diferenciable si y solo si existe un abierto $I$ que contiene a $A$ y una función diferenciable $\widetilde{F}\colon I\to W$ tal que $\widetilde{F}|_A=F$.

\end{bookremark}
\begin{booktheorem}{Teorema (Regla de la cadena)}
Sean $V$, $W$ y $X$ espacios vectoriales reales normados de dimensión finita. Sean $I\subseteq V$ y $\widetilde{I}\subseteq W$ abiertos, y sean

\[
F \colon I \to \widetilde{I},
\qquad
G \colon \widetilde{I} \to X.
\]
Si $F$ es diferenciable en $p\in I$ y $G$ es diferenciable en $F(p)\in \widetilde{I}$, entonces $G\circ F$ es diferenciable en $p$ y

\[
\mathrm{d}(G\circ F)_p
=
\mathrm{d} G_{F(p)} \circ \mathrm{d} F_p.
\]
\end{booktheorem}
\begin{bookproof}{}
Escribamos

\[
q=F(p),\qquad T=\mathrm{d} F_p,\qquad Q=\mathrm{d} G_q.
\]
Queremos probar que

\[
\lim_{v\to 0}
\frac{\|G(F(p+v))-G(F(p))-Q(Tv)\|}{\|v\|}=0.
\]
Definimos

\[
w = F(p+v)-F(p).
\]
Entonces $w\to 0$ cuando $v\to 0$ porque $F$ es continua en $p$. Además,

\[
G(F(p+v))-G(F(p))-Q(Tv)
=
\bigl(G(q+w)-G(q)-Qw\bigr)+Q(w-Tv).
\]
Por el lema de acotación aplicado a $Q$, existe $\mu\ge 0$ tal que

\[
\|Q(w-Tv)\|\le \mu \|w-Tv\|.
\]
Por lo tanto,

\[
\frac{\|G(F(p+v))-G(F(p))-Q(Tv)\|}{\|v\|}
\le
\frac{\|G(q+w)-G(q)-Qw\|}{\|v\|}
+
\mu \frac{\|w-Tv\|}{\|v\|}.
\]
Si $w=0$, el primer sumando es cero. Si $w\neq 0$, escribimos

\[
\frac{\|G(q+w)-G(q)-Qw\|}{\|v\|}
=
\frac{\|G(q+w)-G(q)-Qw\|}{\|w\|}
\cdot
\frac{\|w\|}{\|v\|}.
\]
Además,

\[
\frac{\|w\|}{\|v\|}
\le
\frac{\|F(p+v)-F(p)-Tv\|}{\|v\|}
+
\frac{\|Tv\|}{\|v\|}
\le
\frac{\|F(p+v)-F(p)-Tv\|}{\|v\|}+\lambda
\]
para alguna constante $\lambda\ge 0$, nuevamente por acotación de $T$. El primer factor tiende a $0$ porque $G$ es diferenciable en $q$, y el segundo está acotado cerca de $v=0$. Por otro lado,

\[
\mu \frac{\|w-Tv\|}{\|v\|}
=
\mu \frac{\|F(p+v)-F(p)-Tv\|}{\|v\|}
\longrightarrow 0
\qquad (v\to 0)
\]
porque $F$ es diferenciable en $p$. Concluimos que el cociente completo tiende a $0$.

\bookqed
\end{bookproof}
\section{Derivadas parciales}
\begin{bookremark}{}
Nos enfocamos ahora en el caso $V=\mathbb{R}^n$ y $W=\mathbb{R}^m$ con sus normas usuales. Sea

\[
F \colon U \subseteq \mathbb{R}^n \to \mathbb{R}^m
\]
diferenciable en $p\in U$. En $\mathbb{R}^n$ consideramos la base canónica

\[
\{e_1,\dots,e_n\},\qquad e_i=(0,\dots,0,1,0,\dots,0).
\]
La matriz de $(\mathrm{d} F)_p$ respecto de las bases canónicas de $\mathbb{R}^n$ y $\mathbb{R}^m$ se llama \emph{matriz jacobiana} de $F$ en $p$ y se denota por

\[
JF(p)=\bigl((\mathrm{d} F)_p e_1 \ \cdots \ (\mathrm{d} F)_p e_n\bigr).
\]
Cada columna es el vector $(\mathrm{d} F)_p e_i$, y por la proposición inicial,

\[
(\mathrm{d} F)_p e_i=\lim_{t\to 0}\frac{F(p+t e_i)-F(p)}{t}.
\]
Por otro lado podemos escribir la imagen de $dF_{p}$ en base canonica tenemos

\[
(\mathrm{d} F)_p e_i = \sum_{j=1}^m t_j e_j,
\]
entonces las coordenadas $t_j$ satisfacen

\[
t_j=\left\langle (\mathrm{d} F)_p e_i,e_j\right\rangle=\lim_{t\to 0}\left\langle \frac{F(p+t e_i)-F(p)}{t},e_j\right\rangle.
\]
Si $F=(F_1,\dots,F_m)$, esto equivale a

\[
t_j=\lim_{t\to 0}\frac{F_j(p+t e_i)-F_j(p)}{t}.
\]
\end{bookremark}
\begin{bookdefinition}{Derivada parcial escalar}
Dada $f \colon U \subseteq \mathbb{R}^n \to \mathbb{R}$, con $U$ abierto, la \emph{$i$-ésima derivada parcial} de $f$ en $p\in U$ se define por

\[
\frac{\partial f}{\partial x_i}(p):= \lim_{t\to 0}\frac{f(p+t e_i)-f(p)}{t},
\]
cuando este límite existe.

\end{bookdefinition}
\begin{bookdefinition}{Derivada parcial vectorial}
Si

\[
F=(F_1,\dots,F_m)\colon U \subseteq \mathbb{R}^n \to \mathbb{R}^m,
\]
definimos

\[
\frac{\partial F}{\partial x_i}(p)
:=
\begin{pmatrix}
\dfrac{\partial F_1}{\partial x_i}(p)\\[0.6em]
\vdots\\[0.2em]
\dfrac{\partial F_m}{\partial x_i}(p)
\end{pmatrix},
\]
cuando cada una de esas derivadas parciales existe.

\end{bookdefinition}
\begin{bookremark}{Observación}
Si $F$ es diferenciable en $p$, entonces la matriz jacobiana está dada por

\[
JF(p)=\bigl(a_{jk}\bigr)
=
\left(\frac{\partial F_j}{\partial x_k}(p)\right),
\]
es decir, la $i$-ésima columna de $JF(p)$ es el vector

\[
\frac{\partial F}{\partial x_i}(p).
\]
\end{bookremark}
\begin{bookdefinition}{Función de clase $C^1$}
Sea $F \colon U \subseteq \mathbb{R}^n \to \mathbb{R}^m$, con $U$ abierto. Si todas las derivadas parciales de las componentes de $F$ existen en todo punto de $U$ y las funciones

\[
\frac{\partial F_i}{\partial x_j}\colon U \subseteq \mathbb{R}^{n} \to \mathbb{R}
\]
son continuas, entonces se dice que $F$ es una \emph{función de clase $C^1$}, o bien que es \emph{continuamente diferenciable} en $U$.

\end{bookdefinition}
\begin{bookremark}{Observación}
Si $F$ es diferenciable en $p$, entonces para todo $v\in \mathbb{R}^n$ existe la derivada direccional

\[
(\mathrm{d} F)_p(v)
=
\lim_{t\to 0}\frac{F(p+tv)-F(p)}{t}.
\]
En particular, existen todas las derivadas parciales y $F$ es continua en $p$.

\end{bookremark}
\begin{bookremark}{Observación}
La existencia de todas las derivadas direccionales en un punto no garantiza que $F$ sea diferenciable en ese punto.

\end{bookremark}
\begin{bookexercise}{}
Sea $F:\mathbb{R}^2\to \mathbb{R}$ dada por

\[
f(x,y)= \begin{cases} \dfrac{x^2y}{x^4+y^2}, & (x,y)\neq (0,0),\\[0.6em] 0, & (x,y)=(0,0). \end{cases}
\]
Mostrar que todas las derivadas direccionales existen en $p=(0,0)$, pero $F$ no es continua en $(0,0)$. ¿Qué puede decir de la función?

\end{bookexercise}
\begin{bookproof}{}
\begin{enumerate}
\item Sea $v=(a,b)\in \mathbb{R}^2$. La derivada direccional de $f$ en $(0,0)$ en dirección $v$ es
\end{enumerate}
\[
D_vf(0,0)=\lim_{t\to 0}\frac{f(ta,tb)-f(0,0)}{t}.
\]
\begin{enumerate}[start=2]
\item Como $f(0,0)=0$, para $t\neq 0$ tenemos
\end{enumerate}
\[
f(ta,tb)=\frac{(ta)^2(tb)}{(ta)^4+(tb)^2}=\frac{t^3a^2b}{t^4a^4+t^2b^2}=\frac{ta^2b}{t^2a^4+b^2}.
\]
\begin{enumerate}[start=3]
\item Por lo tanto,
\end{enumerate}
\[
\frac{f(ta,tb)}{t}=\frac{a^2b}{t^2a^4+b^2}.
\]
\begin{enumerate}[start=4]
\item Si $b\neq 0$, al pasar al límite resulta
\end{enumerate}
\[
D_vf(0,0)=\lim_{t\to 0}\frac{a^2b}{t^2a^4+b^2}=\frac{a^2}{b}.
\]
\begin{enumerate}[start=5]
\item Si $b=0$, entonces $f(ta,0)=0$ para todo $t$, y por ende
\end{enumerate}
\[
D_{(a,0)}f(0,0)=0.
\]
\begin{enumerate}[start=6]
\item Luego, todas las derivadas direccionales de $f$ en $(0,0)$ existen.
\item Veamos ahora que $f$ no es continua en $(0,0)$. Si tomamos la curva $y=x^2$, para $x\neq 0$ se obtiene
\end{enumerate}
\[
f(x,x^2)=\frac{x^2\,x^2}{x^4+x^4}=\frac{x^4}{2x^4}=\frac{1}{2}.
\]
\begin{enumerate}[start=8]
\item Así,
\end{enumerate}
\[
\lim_{x\to 0}f(x,x^2)=\frac{1}{2}\neq 0=f(0,0).
\]
\begin{enumerate}[start=9]
\item En consecuencia, $f$ no es continua en $(0,0)$.
\item Por lo tanto, $f$ no es diferenciable en $(0,0)$. Este ejemplo muestra que la existencia de todas las derivadas direccionales en un punto no implica ni continuidad ni diferenciabilidad en ese punto.
\end{enumerate}
\bookqed
\end{bookproof}
\begin{booktheorem}{Diferenciabilidad en $C^1$}
Sea $U\subseteq \mathbb{R}^n$ un abierto. Si $F:U\subseteq \mathbb{R}^n\to \mathbb{R}^m$ es de clase $C^1$ sobre $U$ (esto es, las funciones $\dfrac{\partial F_j}{\partial x_i}:U\subseteq \mathbb{R}^n\to \mathbb{R}$ existen y son continuas), entonces $F$ es diferenciable sobre todo $U$.

\end{booktheorem}
\begin{bookproof}{}
Repasarlo. Cualquier libro clásico de análisis (Rudin, Spivak, Apostol, ...).

\bookqed
\end{bookproof}
\begin{bookdefinition}{Derivadas parciales de orden superior}
Sea $U$ un abierto de $\mathbb{R}^n$ y

\[
F:U\subseteq \mathbb{R}^n \longrightarrow \mathbb{R}^m,\qquad
x\longmapsto (F_1,\dots,F_m),
\]
una función tal que existen todas sus derivadas parciales sobre $U$.

Es decir, podemos definir funciones

\[
\frac{\partial F_j}{\partial x_i}:U\subseteq \mathbb{R}^n \longrightarrow \mathbb{R},
\qquad
p=(x_1,\dots,x_n)\longmapsto \frac{\partial F_j}{\partial x_i}(p).
\]
En este caso, podemos derivar estas funciones para obtener derivadas parciales de segundo orden:

\[
\frac{\partial^2 F_j}{\partial x_k\,\partial x_i}
:=
\frac{\partial}{\partial x_k}\left(\frac{\partial F_j}{\partial x_i}\right),
\]
cuando tales derivadas existan.

Continuando de forma recursiva, se definen las derivadas parciales de orden superior. Las derivadas parciales de orden $k$ son las derivadas parciales de orden $k-1$, cuando estas existen.

\end{bookdefinition}
\begin{bookdefinition}{Clases de diferenciabilidad}
Sea $U\subseteq \mathbb{R}^n$ un abierto y $k\in \mathbb{N}\cup\{0\}$. Una función $F:U\subseteq \mathbb{R}^n\to \mathbb{R}^m$ se dice de clase $C^k$, o diferenciable continuamente $k$ veces, si todas las derivadas parciales de $F$ de orden menor o igual a $k$ existen y son funciones continuas en $U$.

Así, la clase $C^0$ son solo las funciones continuas sobre $U$.

Una función que es de clase $C^k$ para todo $k\in \mathbb{N}\cup\{0\}$ se dice de clase $C^\infty$, o \emph{suave}, o infinitamente diferenciable.

\end{bookdefinition}
\begin{booktheorem}{Schwarz-Clairaut de las derivadas cruzadas}
Si $U$ es un abierto de $\mathbb{R}^n$ y $F:U\subseteq \mathbb{R}^n\to \mathbb{R}^m$ es una función de clase $C^2$ sobre $U$, entonces las derivadas parciales mixtas de segundo orden no dependen del orden de diferenciación:

\[
\frac{\partial^2 F_j}{\partial x_i\,\partial x_k}
=
\frac{\partial^2 F_j}{\partial x_k\,\partial x_i}.
\]
\end{booktheorem}
\begin{bookcorollary}{Corolario}
Si $F:U\subseteq \mathbb{R}^n\to \mathbb{R}^m$ es suave, entonces las derivadas parciales mixtas de $F$ de cualquier orden son independientes del orden de diferenciación.

\end{bookcorollary}
\begin{bookproof}{}
Repasarlo. Cualquier libro clásico de análisis (Rudin, Spivak, Apostol, ...).

\bookqed
\end{bookproof}
\begin{bookremark}{Observación}
Desde el punto de vista de Teoría de Lie, puede entenderse este teorema como que ciertos grupos de transformaciones conmutan entre sí (en este caso, dos translaciones conmutan entre sí; las translaciones "generan" las derivadas parciales).

\end{bookremark}
